% !TEX root = template.tex
\section{Related Work}
\label{sec:related_work}

This section contextualizes our approach within the existing literature. Specifically, the following subsections review the state of the art in WiFi-based activity recognition and explore how Resnet has been adapted to these type of task.

\subsection{HAR Systems}
In [4], the authors propose SHARP, a framework designed to perform the HAR task using commercial IEEE 802.11 (Wi-Fi) devices. This approach demonstrates that fine-grained human actions can be identified with high accuracy, even when performed by different participants within a controlled environment. Architecturally, the SHARP model relies on a sequence of stacked convolutional layers optimized for localized feature extraction to execute classification. While effective, convolutional neural networks are fundamentally constrained by their localized receptive fields, which limit their capacity to model long-range temporal correlations. To overcome this limitation, our work departs from traditional convolutional topologies in favor of a Resnet architecture; the deep layer topology of ResNet allows the network to learn hierarchically abstract spatial and temporal features automatically directly from raw data, overcoming the limitations and overlap of hand-crafted features.

\subsection{Resnet-architecture applied to HAR Systems}
About the Resnet architecture and its application on HAR systems, the literature looks confident for this kind of architecture; for example [8] shows that a Resnet suit well on HAR task rather than a simple CNN. Moreover the authors of [8] focus on the fact that by increasing the depth of the network the performances are improved in terms of precision and recall. Based on these considerations, we have implemented a Resnet, according to the architecture proposed by [1].